\section{Introduction}
In this paper, we consider numerical methods for solving a linear system in the form $AB \bm{x} = \bm{b}$.
This problem is a special case for $A \bm{x} = \bm{b}$, thus we consider three types of methods, direct methods, stationary iterative methods, and Krylov subspace methods.
Each method should be specially treated to make best use of the matrix product structure $AB$.

To apply direct methods, the linear equation ${L}_{A} {U}_{A} {L}_{B} {U}_{B} \bm{x} = \bm{b}$ is exploited
after decomposing each of $A , B$ into lower triangular matrix ${L}$ and upper triangular matrix ${U}$.
After solving four triangular systems, the solution $\bm{x}$ is acquired.
To apply stationary iterative methods, we first solve the linear system $A \bm{y} = \bm{b}$, then solve the linear system $B \bm{x} = \tilde{\bm{y}}$ using the aomputed approximation $\tilde{\bm{y}}$.
To apply Krylov subspace methods, we have two different approaches. The first approach is similar to the approach when applying stationaty iterative methods; The second approach computes at each iterate the matrix-vector product $AB\bm{p}$ in the order $A (B\bm{p})$.
By now, we have two approaches for solving a linear system $AB \bm{x} = \bm{b}$. The separate method. The combined method. The combined method can only be implemented when Krylov subspave methods are applied.

We take special interest in the following case.

And these phenomenon is not analysed into details, thus we cannot arrive a conclusion at which method should be used at which circumstance.

However what is the main motivation? I guess it is still not clear. The main motivation. It is not "because we need a method" but should be "somehow what problem attracts our attention".

It is a simple idea thus what about, an example that separate method works extremely slow but the other one is very fast?

It is not about structural but how the operates are affected accross themselves.

